\subsection{The two possible full images at degree \(163\)}
Write \(\mathcal R=R_{\rm global}\) and
\(\mathcal V=V_{\rm global}\), let \(P=P_{163}\), and put
\(A=Q_{163}P\).  Thus \(P/A\) is the two-dimensional irreducible
\(\mathbf F_2[C_3]\)-module.  The cyclic group \(C_3\) is the quotient
of \(\operatorname{Gal}(\mathbf Q(\zeta_{163})/\mathbf Q)\) through
which this module is acted on.

\begin{proposition}\label{prop:cm163-full-image-dichotomy}
The full images \(h_{163}(\mathcal R)\) and
\(h_{163}(\mathcal V)\) coincide.  There are precisely two
possibilities.
\begin{enumerate}
\item Their common image is \(A\).  The groups
\(\mathcal R'=\mathcal R\cap\ker h_{163}\) and
\(\mathcal V'=\mathcal V\cap\ker h_{163}\) then give a further
arithmetic normalization \(G=\mathcal R'H\).  The full relative
abelian quotient on \(\mathcal R'\), its maximal pro-\(2\) quotient
on \(\mathcal V'\), and the exact original quadratic restriction
kernel and image are unchanged.
\item Their common image is \(P\).  There is a continuous
surjection \(G\longrightarrow A_4\) whose restriction to \(H\)
has image a subgroup of order three.
\end{enumerate}
In the second case the displayed finite quotient has zero relative
quadratic character space.
\end{proposition}
\begin{proof}
Since \(\mathcal R/\mathcal V\) is pro-odd, the quotient
\(h_{163}(\mathcal R)/h_{163}(\mathcal V)\) is both pro-odd and
pro-\(2\), and is trivial.  The exact arithmetic image is \(A\).
The common full image modulo \(A\) is an invariant subspace of
the irreducible module \(C=P/A\), and hence is zero or \(C\).

If the common image is \(A\), its equality with the arithmetic
image on both arithmetic intersections proves
\(\mathcal R=\mathcal R'(\mathcal V\cap H)\), and
therefore \(G=\mathcal R'H\).  The relative gluing and pro-odd
comparison arguments apply as before.  The exact original
restriction statements follow from the decomposition into the
relative character space and the fixed lift space \(\mathcal S\).

Suppose instead that the common full image is \(P\), and let
\(t:\mathcal R\to C\) be \(h_{163}\) followed by reduction modulo
\(A\).  It is surjective and vanishes on \(\mathcal R\cap H\).
Write \(\gamma:G\to C_3\) for the cyclotomic quotient acting on
\(C\).  For \(g=rh\), \(r\in\mathcal R\), \(h\in H\), set
\(z(g)=t(r)\).  This is well defined because \(t\) vanishes on
the overlap.  It is continuous: the continuous function
\((r,h)\mapsto t(r)\) on the compact space \(\mathcal R\times H\)
descends through its quotient map onto \(G\).  Equivariance and
the multiplication rule for \(rh\) give
\[
 z(g_1g_2)=z(g_1)+\gamma(g_1)z(g_2).
\]
Here \(\gamma(\mathcal R)=1\), as \(\mathcal R\subset\ker\lambda\).
It follows that
\[
 g\longmapsto (z(g),\gamma(g))
\]
is a continuous homomorphism to \(C\rtimes C_3\).  Its image
contains all translations \(C\) and the entire complement \(C_3\);
the latter follows from \(G=\mathcal RH\) and
\(\gamma(\mathcal R)=1\).  It is therefore surjective.
The faithful irreducible action identifies
\(C\rtimes C_3\) with \(A_4\), for example through its affine
action on the four elements of \(C\).  The image of \(H\) is
exactly the complement.

Finally \((C_3-1)C=C\), since the irreducible nontrivial
\(C_3\)-module \(C\) has no trivial constituent.  Every invariant
homomorphism \(C\to\mathbf F_2\) is consequently zero.  This
proves the last assertion without deciding the original relative
quadratic character space.
\end{proof}
